\section{The conductor-one arithmetic packet}\label{sec:sr-arithmetic}
The arithmetic specialization in this section follows the contributors'
\texttt{PrincipalHostField}, \texttt{PrincipalCohomology},
\texttt{LengthBudget}, \texttt{PrincipalMinus},
\texttt{PrincipalPlus}, and \texttt{PowerWindow} files.
We prove the elementary field, class-number, descent, and numerical
comparisons here. The remaining uniform height and rank inputs are
distinguished from these consequences at the end of the section.
Classical background used here consists of finite Galois theory,
the Dirichlet unit theorem, Hilbert theorem 90, the Minkowski ideal-class
bound, and the unramified Artin reciprocity theorem.
None is supplied by a reference to a Lean module.

\subsection{An explicit field and its degree}
Let \(q\equiv3\pmod8,d\equiv7\pmod{16}\) be primes and \(e\geq0\).
They are distinct odd primes. Set \(r=e+2\), \(N=2^rqd\), and
\(L=\mathbb Q(\zeta_N)\). Let \(A\) be the subgroup of
\((\mathbb Z/N)^\times=\operatorname{Gal}(L/\mathbb Q)\) defined by
\[
 a\bmod2^r\in\{1,-1\},\qquad
 \left(\frac aq\right)=\left(\frac ad\right)=1,
\]
and put \(E=L^A\).
The Chinese remainder theorem and surjectivity of each quadratic character
give
\[
 \operatorname{Gal}(E/\mathbb Q)\simeq
 ((\mathbb Z/2^r)^\times/\{\pm1\})\times(\mathbb Z/2)^2,\qquad
 [E:\mathbb Q]=2^r=2^{e+2}.
\]
Indeed the first factor has order \(2^{r-2}\), including \(r=2\),
and each of the other factors has order two.
In particular no odd prime divides the degree.

Complex conjugation is the class of \(-1\).
It is trivial on the first factor and nontrivial on both quadratic
factors, since \(q,d\equiv3\pmod4\).
It is therefore a nonidentity central involution.
Its fixed field \(E^+\) lies in \(\mathbb R\), and
\([E:E^+]=2\). Every embedding of the Galois field \(E\) has the same
nontrivial conjugation action. Thus \(E\) is a CM field.
The least positive conductor \(f_E\) for a cyclotomic embedding
satisfies \(f_E\leq N\).

This construction is independent of the ambient cyclotomic field used
to describe it. In a larger cyclotomic field of conductor divisible by
\(N\), the subgroup specified by these same three tests is the inverse
image of \(A\). Restriction of automorphisms is surjective, and Galois
correspondence identifies its fixed field with \(E\).
Changing the chosen primitive root changes an exponent description by
multiplication by a unit and leaves the automorphism's exponent itself
unchanged. Therefore the fixed subgroup, and hence its fixed field,
are unchanged. This proves the canonical-field reduction without
retaining any choice of a realization.

\subsection{Strict lengths and a single integer}
For a prime \(\ell\), write \(S_\ell\) for the primes of \(E\) above \(\ell\).
For \(a\geq1\), define the strict group
\[
 T_{\ell,a}(E)=\ker\left[
 H^1(G_{E,S_\ell},(\mathbb Q_\ell/\mathbb Z_\ell)(-a))
 \longrightarrow
 \prod_{v\mid\ell}H^1(E_v,(\mathbb Q_\ell/\mathbb Z_\ell)(-a))\right].
\]
Cohomology is continuous, and the local maps are restrictions after
inflation from the Galois group unramified outside \(S_\ell\).
Let \(s_{\ell,E}(a)\) be its \(\mathbb Z_\ell\)-length, with value infinity
allowed until finiteness has been proved.
The group and length are independent of choices of closures and primes
above local places: field isomorphisms induce isomorphisms of global
and local groups, commute with restriction, and preserve these kernels.

Here is a useful exact reformulation, following \texttt{LengthBudget}.
For nonnegative integers \(L_\ell\) and a real \(B\),
\begin{equation}\label{eq:sr-integerbudget}
 \left(\forall S\ {\rm finite},\ \sum_{\ell\in S}L_\ell\log\ell\leq B\right)
 \quad\Longleftrightarrow\quad
 \left(\exists M\in\mathbb Z_{\geq1},\
 \log M\leq B,\ \ell^{L_\ell}\mid M\ {\rm for\ every}\ \ell\right).
\end{equation}
For the forward implication, the empty set first gives \(B\geq0\).
If \(L_\ell>0\), the singleton inequality gives \(\ell\leq e^B\).
There are only finitely many such primes. Take \(M\) to be the product
of \(\ell^{L_\ell}\) over this finite support; its logarithm is at most
\(B\) by the hypothesis. Conversely, distinct prime powers are coprime,
so each finite product divides \(M\), and its logarithm is at most
\(\log M\leq B\).
Apply this with \(L_\ell=\sum_{1\leq a<t}s_{\ell,E}(a)\), after proving
every length in that sum finite.
Finiteness cannot be dropped: replacing infinity by zero is not a
mathematical interpretation of a length.

\subsection{The one-ramified-prime class-number lemma}
\begin{srlemma}\label{lem:sr-oneprime}
Let \(L/K\) be a cyclic extension of degree \(p\), with \(p\) prime,
and both fields totally imaginary. Suppose exactly one finite prime
of \(K\) ramifies, and it is totally ramified. If
\(p\nmid h_K\), then \(p\nmid h_L\).
\end{srlemma}
\begin{proof}
Write \(G=\operatorname{Gal}(L/K)\), \(U_L\) for the units,
\(I_L\) for fractional ideals, and \(P_L\) for principal fractional ideals.
The group \(I_L\) is a direct sum of permutation modules on the primes.
Its first cohomology is zero: for each orbit this reduces to
\(H^1(J,\mathbb Z)=\operatorname{Hom}(J,\mathbb Z)=0\), with \(J\) finite.
Taking invariants in \(0\to P_L\to I_L\to\operatorname{Cl}_L\to0\)
therefore gives
\[
 |\operatorname{Cl}_L^G|
   =|I_L^G/P_L^G|\ |H^1(G,P_L)|.
\]
Extension of ideals identifies \(I_K\) with a subgroup of \(I_L^G\).
Invariants on each prime orbit show
\([I_L^G:I_K]=\prod_{v\ {\rm finite}}e_v=p\).
The sequence \(1\to U_L\to L^\times\to P_L\to1\), together with
Hilbert theorem 90, gives
\(P_L^G/P_K\simeq H^1(G,U_L)\). Thus
\[
 |I_L^G/P_L^G|=h_Kp/|H^1(G,U_L)|.
\]
In the same sequence,
\[
 H^1(G,P_L)\simeq
 \ker\bigl(H^2(G,U_L)\to H^2(G,L^\times)\bigr)
 = (U_K\cap N_{L/K}L^\times)/N_{L/K}U_L.
\]
For a cyclic group, \(H^2(G,U_L)=U_K/N_{L/K}U_L\).
Consequently
\[
 |\operatorname{Cl}_L^G|
 =\frac{h_Kp}{[U_K:U_K\cap N_{L/K}L^\times]}\,
       \frac{|H^2(G,U_L)|}{|H^1(G,U_L)|}.
\]

For clarity, the last ratio is \(1/p\).
Finite modules have Herbrand quotient one, as is seen by counting the
kernel and image of \(g-1\) and of the norm.
The logarithmic unit lattice, modulo its finite torsion, spans the
hyperplane of sum zero in the real vector space on the complex places.
By the Dirichlet unit theorem its rational representation is
\(\mathbb Q[G]^{r_2(K)}-\mathbb Q\).
Herbrand quotient is unchanged under passage to a commensurable
integral lattice, since its finite quotient has Herbrand quotient one.
For \(\mathbb Z[G]\) that quotient is one; for the trivial lattice
\(\mathbb Z\) it is \(p\).
The exact sequence from the sum-zero lattice to
\(\mathbb Z[G]^{r_2(K)}\) and then to \(\mathbb Z\) therefore gives \(1/p\).
These facts also show that the cohomology groups counted above are finite.
We obtain the ambiguous-class formula, in precisely this case,
\[
 |\operatorname{Cl}_L^G|
 =\frac{h_K}{[U_K:U_K\cap N_{L/K}L^\times]}.
\]
Its right side is prime to \(p\).
If the \(p\)-Sylow subgroup \(A\) of \(\operatorname{Cl}_L\) were
nontrivial, the action of \(G\) on \(A\) would have a nontrivial fixed
point. Indeed its orbits have sizes one or \(p\), so
\(|A^G|\equiv|A|\equiv0\pmod p\), while the identity is fixed.
This would make \(p\) divide \(|\operatorname{Cl}_L^G|\), a contradiction.
\end{proof}

The units argument uses the usual rational representation form of the
Dirichlet unit theorem. One can see the rational representation directly
from its real logarithmic embedding: the invariants for every subgroup
have dimension \(r_2(L^J)-1\); these determine the rational character
of the cyclic group and give the stated representation.
Thus the passage between the logarithmic real space and integral
lattices does not assume a rational-valued logarithm.

In \(\mathbb Q(\zeta_{p^{n+1}})/\mathbb Q(\zeta_{p^n})\), the degree
is \(p\), the unique prime above \(p\) is totally ramified, and no other
finite prime ramifies. These facts follow from the Eisenstein
cyclotomic polynomial for \(\zeta_{p^n}-1\) and its discriminant,
which is a power of \(p\).
Starting with \(\mathbb Q(i)\), whose Gaussian integer ring is Euclidean
by nearest-integer division in each coordinate, Lemma~\ref{lem:sr-oneprime}
therefore proves
\[
 h_{\mathbb Q(\zeta_{2^n})}\ \hbox{is odd}\qquad(n\geq2).
\]
This proves the Weber input used in the contributors' dyadic reduction.

The same argument applies to the primes 3 and 7 once the first class
numbers are checked. For \(\mathbb Q(\zeta_3)\), the Minkowski bound is
\((2/\pi)\sqrt3<2\), so every class has an ideal of norm one.
For \(\mathbb Q(\zeta_7)\), the degree is six and the discriminant has
absolute value \(7^5\); its Minkowski bound is
\[
 (4/\pi)^3\frac{6!}{6^6}7^{5/2}<5.
\]
Primes above 2 have residue degree three, and primes above 3 have
residue degree six, because those are the orders of 2 and 3 modulo 7.
Thus no ideal has norm 2, 3 or 4.
Every ideal class has a representative of norm at most four, so the
class number is one.
Here the discriminant is obtained by taking the product of the
derivatives of \(\Phi_p=(X^p-1)/(X-1)\) at its roots;
its absolute value is \(p^{p-2}\).
We have proved
\[
 p\nmid h_{\mathbb Q(\zeta_{p^n})}
 \qquad(p=3,7,\ n\geq1).
\]

\subsection{Strict descent in the regular cyclotomic towers}
\begin{srproposition}\label{prop:sr-regularsha}
For \(p=3,7\), \(F=\mathbb Q(\zeta_{p^n})\), \(n\geq1\), and every
\(a\geq1\), the strict group \(T_{p,a}(F)\) is zero.
\end{srproposition}
\begin{proof}
Let \(D=(\mathbb Q_p/\mathbb Z_p)(-a)\) and represent a strict class
by a continuous cocycle \(z:G_F\to D\), inflated from the group
unramified outside \(p\).
Its image is finite because \(G_F\) is compact and \(D\) is discrete.
Hence all its values belong to \(D[p^b]\) for some \(b\).
At the unique prime above \(p\), strictness says that the cocycle is a
coboundary, \(z(g)=(g-1)v\), on the decomposition group, for some \(v\in D\).
Choose a finite higher cyclotomic layer \(F_m\) in which the cyclotomic
action is trivial on both \(D[p^b]\) and \(v\).

The restriction \(z|_{G_{F_m}}\) is then a homomorphism into \(D[p^b]\).
It is zero on decomposition at \(p\).
At every other finite prime, it is zero on inertia: the coefficient
action is unramified there and the original cocycle was inflated from
the group unramified outside \(p\).
All higher primes above \(p\) are conjugate, so the same vanishing
holds at each choice of a local prime.
Unramified Artin reciprocity makes this homomorphism factor through
\(\operatorname{Cl}_{F_m}\).
The latter group has order prime to \(p\), by the preceding result.
A homomorphism from it into a finite \(p\)-group is zero.
Thus \(z|_{G_{F_m}}=0\).

Let \(F_\infty=\bigcup_mF_m\).
The action on \(D\) is trivial on \(G_{F_\infty}\), and the preceding
vanishing lets \(z\) descend to a continuous cocycle on
\(\Gamma=\operatorname{Gal}(F_\infty/F)\simeq\mathbb Z_p\).
Take a topological generator \(g\) of \(\Gamma\).
The scalar \(g-1\) on \(D\) is
\(\chi_{\rm cyc}(g)^{-a}-1\ne0\).
Multiplication by every nonzero element of \(\mathbb Z_p\) is surjective
on \(\mathbb Q_p/\mathbb Z_p\). Choose \(v_0\in D\) such that
\((g-1)v_0=z(g)\).
After subtracting this coboundary the cocycle vanishes on \(g\), hence
on all its integral powers by the cocycle identity, and hence on
\(\Gamma\) by continuity and density.
The original class is zero.
\end{proof}
This proves the zero bottom-layer budgets for the two exceptional
auxiliary primes, rather than merely naming an unramified-character input.

\subsection{The odd-prime numerical budget and its remaining premise}
For odd \(\ell\), complex conjugation on \(T_{\ell,a}(E)\) has complementary
idempotents \((1+(-1)^ac)/2\) and \((1-(-1)^ac)/2\).
Call their lengths \(s^+_{\ell,E}(a)\) and \(s^-_{\ell,E}(a)\);
the minus eigenvalue is \((-1)^{a+1}\).
When finite, they sum to \(s_{\ell,E}(a)\).

Define \(b(f,k)=16f^{k+1}k!\).
The precise prime-to-degree estimate used in the contributors' odd
packet is the following assertion, here denoted \(\mathrm{PB}(E)\):
all the minus lengths are finite, and, for every \(a\geq1\) and finite
set \(S\) of odd primes not dividing \([E:\mathbb Q]\),
\begin{equation}\label{eq:sr-PB}
 \sum_{\ell\in S}s^-_{\ell,E}(a)\log\ell
       \leq[E:\mathbb Q]\log b(f_E,a+1).
\end{equation}
For the canonical conductor-one field, every odd prime is eligible.
Assuming \eqref{eq:sr-PB}, for \(t\geq2\) we get the fully explicit bound
\begin{equation}\label{eq:sr-minus}
 \sum_{\substack{\ell\in S\\\ell\ne2}}\log\ell
       \sum_{a=1}^{t-1}s^-_{\ell,E}(a)
 \leq3[E:\mathbb Q]t^2\log(tf_E).
\end{equation}
To check the entire numerical estimate, set \(u=tf_E\geq2\).
For \(a<t\), \(16\leq u^4\), \(f_E^{a+2}\leq u^{a+2}\), and
\((a+1)!\leq(a+1)^{a+1}\leq u^{a+1}\).
Thus \(b(f_E,a+1)\leq u^{2a+7}\).
Consequently
\[
 \sum_{a=1}^{t-1}\log b(f_E,a+1)
 \leq(t-1)(2t+5)\log u\leq3t^2\log u.
\]
The last inequality follows from
\(3t^2-(t-1)(2t+5)=t^2-3t+5>0\).
Interchange the two finite sums in \eqref{eq:sr-PB} to obtain
\eqref{eq:sr-minus}.
This proves the polynomial budget from \(\mathrm{PB}(E)\);
it does not prove \(\mathrm{PB}(E)\).
The characterwise Selmer estimate behind that premise must still be
verified in a completed standalone argument.

\subsection{Selecting a whole interval away from finitely many exceptions}
Let \(s,B\geq1\), \(U>0\), and let \(A\) be a set of at most \(R\)
real numbers. There exists \(y\) such that
\[
 \frac{U}{s(sB)^R}\leq y,\qquad sy\leq U,
\]
and, for \(y\leq u\leq sy\) and \(a\in A\),
\begin{equation}\label{eq:sr-logavoid}
 a\leq u\quad\Longrightarrow\quad Ba\leq u.
\end{equation}
Here is the complete induction on \(R\).
For \(R=0\), use \(y=U/s\).
If some \(b\in A\) lies in \((U/(sB),U]\), remove it and apply the
induction hypothesis to \(U'=U/(sB)\) and \(A\setminus\{b\}\).
The resulting lower bound is exactly \(U/[s(sB)^R]\), and
\(sy\leq U'<b\), so \(b\leq u\) never occurs on the resulting interval.
The other points satisfy \eqref{eq:sr-logavoid} by induction.
If there is no such \(b\), choose \(y=U/s\).
Any \(a\leq u\leq U\) must satisfy \(a\leq U/(sB)\), hence
\(Ba\leq U/s=y\leq u\).
The required lower bound follows from \(sB\geq1\).

Now let \(0<\eta\leq C_t\), \(X>1\), and let \(F\) be a set of at most
\(R\) nonnegative integers.
Apply the result with
\(U=\log X,B=C_t/\eta\), and
\(A=\{\log a/C_t:a\in F,\ a>0\}\).
Exponentiating gives \(Y\) with
\[
 X^{1/J}\leq Y,\qquad Y^s\leq X,\qquad
 J=s(sC_t/\eta)^R,
\]
such that for every \(Y\leq x\leq Y^s\), every \(a\in F\) with
\(a\leq x^{C_t}\) already satisfies \(a\leq x^\eta\).
The value \(a=0\) satisfies this automatically.
For \(s=2,R=24,\eta=\epsilon/(2C_K)\), this is exactly the explicit
interval exponent
\begin{equation}\label{eq:sr-J}
 J=2^{25}(2C_KC_t/\epsilon)^{24}.
\end{equation}
The exceptional set must be fixed before \(x\) varies on the whole
interval. A new exceptional set at each \(x\) would not justify this argument.

\subsection{Absorbing the conductor-one budgets uniformly}
We finish the numerical passage needed after the difficult bounds have
been obtained. Fix \(\epsilon>0,C\geq0\), take
\(t=\lfloor x^{401}\rfloor\), and suppose
\[
 x^{2+8\epsilon}\leq2^e\leq2x^{2+8\epsilon},
 \qquad \log(qd)\leq x^{\epsilon/4}.
\]
For the explicit field,
\[
 [E:\mathbb Q]\leq8x^{2+8\epsilon},\qquad
 f_E\leq N=2^{e+2}qd,\qquad
 \log(Nt)\leq(403+8\epsilon)\log x+\log8+x^{\epsilon/4}.
\]
In particular \(\log(Nt)\leq2x^{\epsilon/4}\) for all sufficiently large
\(x\), with a threshold independent of \(q,d,e\) satisfying these inequalities.

Assume the dyadic strict length sum is at most
\(B_2=\exp(x^{4\epsilon/5})\), the odd minus sum is bounded by
\eqref{eq:sr-minus}, and the odd plus sum is at most \((Nt)^C\).
For every finite prime set, the total strict sum is then at most
\begin{equation}\label{eq:sr-budgetsum}
 (\log2)e^{x^{4\epsilon/5}}
 +3[E:\mathbb Q]t^2\log(tf_E)+(Nt)^C.
\end{equation}
The degree and conductor estimates bound its middle term, eventually, by
\(48x^{804+8\epsilon+\epsilon/4}\).
Its last term is at most \(\exp(2Cx^{\epsilon/4})\).
Each of the three displayed terms divided by
\(\exp(x^\epsilon/2)\) tends to zero:
take logarithms, divide by \(x^\epsilon\), and use respectively
\(x^{-\epsilon/5}\to0\), \((\log x)/x^\epsilon\to0\), and
\(x^{-3\epsilon/4}\to0\).
Thus \eqref{eq:sr-budgetsum} is at most \(\exp(x^\epsilon/2)\)
for all sufficiently large \(x\), uniformly in the auxiliary data.
Equation \eqref{eq:sr-integerbudget} supplies a common integer \(M\)
with \(\log M\leq\exp(x^\epsilon/2)\), if the stipulated finiteness
and the three length estimates have been proved.

\paragraph{What is still required here.}
The class-number and strict descent arguments above prove the auxiliary
regularity and zero bottom-layer inputs.
They do not prove the uniform plus estimate \((Nt)^C\) furnished by
the archive's pointwise height theorem.
The odd estimate \eqref{eq:sr-PB}, the quantitative dyadic rank bound
that produces at most 24 exceptional values, and the cohomological
size estimate called Lemma \(K\) in the source also need their full
standalone proofs before this packet constructs the required window.
The exceptional-interval selection and final absorption are proved here
conditional on those precise inputs.
